\chapter{Executive Summary}

MathDevMCP is proposed as an internal mathematical development copilot platform for groups that build mathematically intensive software and must simultaneously produce large-scale technical documentation, proposals, and monographs under severe time pressure. The immediate motivation comes from recurring experience in DSGE development, macrofinance research, and adjacent quantitative workflows where the same team must maintain consistency across code, equations, assumptions, algorithm descriptions, and long-form exposition. Existing LLM-assisted editors have already proven useful for drafting and local editing, but they remain unreliable in the exact areas that matter most for this group: long-range dependency tracking, cross-file mathematical consistency, derivation checking, code--document traceability, and disciplined auditing of claims.

The proposal recommends building MathDevMCP as a small but rigorous internal platform rather than as a generic chatbot enhancement. The core design combines three elements: a Python implementation layer for extraction, indexing, verification, and consistency logic; a thin MCP server that exposes these capabilities to Claude Code and related tooling; and a workflow layer focused on the actual group tasks of monograph drafting, proposal writing, code-document consistency review, documentation auditing, and mathematically informed code generation. This approach is intentionally conservative. It does not assume autonomous theorem proving, fully automatic proof repair, or end-to-end generation of correct monographs. Instead, it aims to build a trustworthy system that makes mathematically heavy work faster, more reproducible, and less error-prone.

The expected value for the group is broad. For long monographs and proposals, MathDevMCP should improve retrieval, grounding, and revision consistency across hundreds of pages. For mathematical codebases, it should improve the ability to connect equations and assumptions in documents to concrete implementations and test artifacts. For auditing and review, it should support symbolic verification, numeric falsification, gap localization, and consistency checks across code and prose. For code generation, it should constrain generation against mathematical specifications and benchmark the results against executable acceptance criteria. The intent is not merely to generate nicer text, but to reduce the incidence of mathematically confident but incorrect outputs. This design choice follows the broader view that reliable mathematical software work depends on numerical stability, explicit assumptions, and disciplined evidence trails rather than prose fluency alone \citep{higham2002,mackay2003}.

The proposal therefore frames MathDevMCP as a group productivity and reliability investment. Its success will be evaluated not by prose quality alone, but by measurable improvements in retrieval accuracy, code--document consistency detection, audit verdict accuracy, code generation correctness, drafting throughput, and reduction in hallucinated or unsupported mathematical claims.

\begin{proposalbox}
\textbf{Recommendation.} Develop MathDevMCP as a proposal-first and benchmark-driven internal platform. Prioritize document indexing, code--document traceability, symbolic and numeric checks, and structured audit workflows before attempting ambitious autonomous generation features.
\end{proposalbox}
